% BEGIN VOL07-EXPANSION VII16
\section{Supplementary worked examples}
\label{sec:vii16-pedagogy}

\begin{example}[Gauss map of a sphere]\label{ex:vii16-ped-01}
For the sphere of radius \(R\) with outward unit normal, the Gauss map is \(N(p)=p/R\).  Hence \(dN_p=(1/R)\,\mathrm{id}\) on the tangent plane.  With the common convention \(S=-dN\), the shape operator is \(-R^{-1}\mathrm{id}\); with \(S=dN\), the sign reverses.  In either convention every direction is principal.
\end{example}

\begin{example}[Shape operator of a cylinder]\label{ex:vii16-ped-02}
On the cylinder \(X(u,v)=(R\cos u,R\sin u,v)\), the unit normal depends on \(u\) but not on \(v\).  Consequently one principal curvature has magnitude \(1/R\) in the circular direction and the other is \(0\) along the rulings.
\end{example}

\begin{example}[A saddle has opposite principal signs]\label{ex:vii16-ped-03}
For the graph \(z=x^2-y^2\), at the origin the tangent plane is horizontal and the Hessian has eigenvalues \(2\) and \(-2\).  Up to the chapter's normal/sign convention, the shape operator has opposite eigenvalue signs there, anticipating negative Gaussian curvature.
\end{example}

\section{Graded supplementary exercises}
The following exercises add a deliberate progression from concrete calculation to structural proof, hypothesis testing, application, and synthesis.

\subsection*{Standard computations and constructions}

\begin{exercise}[Sphere Gauss map]\label{exr:vii16-ped-01}
Write the Gauss map of the radius-\(R\) sphere with outward normal.
\end{exercise}
\begin{hint}
Normalize the position vector.
\end{hint}
\begin{solution}
\(N(p)=p/R\).
\end{solution}

\begin{exercise}[Cylinder normal]\label{exr:vii16-ped-02}
For \(X(u,v)=(R\cos u,R\sin u,v)\), find a unit normal.
\end{exercise}
\begin{hint}
Normalize \(X_u\times X_v\).
\end{hint}
\begin{solution}
One choice is \(N=(\cos u,\sin u,0)\).
\end{solution}

\begin{exercise}[Cylinder principal directions]\label{exr:vii16-ped-03}
Identify the circular and axial principal directions on a cylinder.
\end{exercise}
\begin{hint}
Differentiate the normal in the \(u\) and \(v\) directions.
\end{hint}
\begin{solution}
The circular direction changes the normal and has curvature magnitude \(1/R\); the axial direction leaves the normal fixed and has curvature \(0\).
\end{solution}

\begin{exercise}[Plane shape operator]\label{exr:vii16-ped-04}
Compute the shape operator of a plane.
\end{exercise}
\begin{hint}
The unit normal is constant.
\end{hint}
\begin{solution}
\(dN=0\), so \(S=0\) in either sign convention.
\end{solution}

\begin{exercise}[Saddle Hessian]\label{exr:vii16-ped-05}
Find the eigenvalues of the Hessian of \(f=x^2-y^2\) at the origin.
\end{exercise}
\begin{hint}
Differentiate twice.
\end{hint}
\begin{solution}
The Hessian is \(\mathrm{diag}(2,-2)\), with eigenvalues \(2,-2\).
\end{solution}

\subsection*{Proofs}

\begin{exercise}[Shape operator self-adjoint]\label{exr:vii16-ped-06}
Prove the shape operator of a surface in \(\mathbb R^3\) is self-adjoint with respect to the first fundamental form.
\end{exercise}
\begin{hint}
Use symmetry of the second fundamental form.
\end{hint}
\begin{solution}
Since \(II(u,v)=I(Su,v)\) and \(II\) is symmetric, \(I(Su,v)=I(u,Sv)\).
\end{solution}

\begin{exercise}[Principal directions are orthogonal]\label{exr:vii16-ped-07}
Prove eigenvectors of \(S\) for distinct eigenvalues are orthogonal.
\end{exercise}
\begin{hint}
Use self-adjointness.
\end{hint}
\begin{solution}
If \(Su=\lambda u\) and \(Sv=\mu v\), then \(\lambda I(u,v)=I(Su,v)=I(u,Sv)=\mu I(u,v)\); if \(\lambda\neq\mu\), then \(I(u,v)=0\).
\end{solution}

\begin{exercise}[Normal reversal]\label{exr:vii16-ped-08}
Show \(N\mapsto-N\) changes \(S\) to \(-S\).
\end{exercise}
\begin{hint}
Differentiate the new Gauss map.
\end{hint}
\begin{solution}
\(d(-N)=-dN\), so either convention for \(S=\pm dN\) changes sign.
\end{solution}

\begin{exercise}[Gauss map differential is tangent]\label{exr:vii16-ped-09}
Show \(dN_p(v)\) lies in \(T_pS\).
\end{exercise}
\begin{hint}
Differentiate \(N\cdot N=1\).
\end{hint}
\begin{solution}
\(2N\cdot dN_p(v)=0\), so \(dN_p(v)\) is orthogonal to \(N\), hence tangent.
\end{solution}

\subsection*{Counterexamples and hypothesis tests}

\begin{exercise}[Gauss map is always injective]\label{exr:vii16-ped-10}
Test on a cylinder.
\end{exercise}
\begin{hint}
Points on one ruling share the same normal.
\end{hint}
\begin{solution}
False.  The normal is independent of the axial coordinate, so infinitely many points have the same Gauss image.
\end{solution}

\begin{exercise}[Shape operator depends only on the metric]\label{exr:vii16-ped-11}
Test with plane and cylinder local isometry.
\end{exercise}
\begin{hint}
They share the same first form but bend differently.
\end{hint}
\begin{solution}
False.  The shape operator is extrinsic and distinguishes a plane from an isometric cylinder.
\end{solution}

\begin{exercise}[Principal curvatures keep their signs under normal reversal]\label{exr:vii16-ped-12}
Test \(N\mapsto-N\).
\end{exercise}
\begin{hint}
All eigenvalues of \(S\) change sign.
\end{hint}
\begin{solution}
False.  Principal curvatures reverse sign, although products such as Gaussian curvature do not.
\end{solution}

\subsection*{Applications and investigations}

\begin{exercise}[Surface normal estimation]\label{exr:vii16-ped-13}
Why is the Gauss map relevant in geometry processing?
\end{exercise}
\begin{hint}
It records a unit normal at every surface point.
\end{hint}
\begin{solution}
Variation of the normal field approximates the differential of the Gauss map and hence the shape operator, providing curvature information from sampled geometry.
\end{solution}

\begin{exercise}[Directional bending]\label{exr:vii16-ped-14}
Interpret an eigenvector of the shape operator.
\end{exercise}
\begin{hint}
Relate it to normal curvature.
\end{hint}
\begin{solution}
A principal direction is a tangent direction in which normal bending is extremal and the shape operator acts by simple scaling.
\end{solution}

\subsection*{Challenge problems}

\begin{exercise}[Sphere is totally umbilic]\label{exr:vii16-ped-15}
Use the Gauss map to show every point of a sphere is umbilic.
\end{exercise}
\begin{hint}
\(dN\) is a scalar multiple of the identity.
\end{hint}
\begin{solution}
Since \(N(p)=p/R\), its differential on tangent spaces is \(R^{-1}\mathrm{id}\); thus \(S\) is scalar and all tangent directions are principal with equal curvature.
\end{solution}

\begin{exercise}[Zero shape operator]\label{exr:vii16-ped-16}
Show that a connected surface patch with identically zero shape operator lies in a plane.
\end{exercise}
\begin{hint}
Then \(dN=0\), so the normal is constant.
\end{hint}
\begin{solution}
A constant normal \(N_0\) implies \(d(X\cdot N_0)=dX\cdot N_0=0\), so \(X\cdot N_0\) is constant on the connected patch; the image lies in an affine plane.
\end{solution}

% END VOL07-EXPANSION VII16
